\textbf{COLUNI 2026 -- Questão 08}\par\medskip O conto \textit{A morte e o Rei} trata da temática da morte. Outros contos do livro \textit{Com certeza tenho amor} também abordam essa simbologia. Sobre isso, leia os excertos a seguir:

\begin{enumerate}
\item É nesse gesto, como se o completasse, que a mulher se achega. Colhe a mão, desfaz a atadura e docemente, muito docemente, começa a lamber a ferida. Sobre o talho que aos poucos se fecha, pelos vermelhos despontam e, como um arrepio que percorre o corpo, se alastram braço acima, tomam o peito, invadem toda a pele. (Conto \textit{Na neve, os caçadores}, p. 88)
\item Mulher de cadeiras largas, sem esforço pariu o primeiro filho na cama que havia sido da sua família. Oferecia-lhe o peito ainda deitada, quando a vizinha veio visitá-la. Debruçou-se elogiando o pequeno, entregou à mãe a laranja que havia trazido, e cedo despediu-se. Mas ainda na porta voltou-se, olhou a cama. – Leito de vida, leito de morte – disse sem alegria. E se foi. (Conto \textit{Rosas na cabeceira}, p. 47)
\item Assim, ninguém viu aquele súbito mover-se entre cinzas, o menear, a cabeça da salamandra erguendo-se. Ninguém viu o braço, o ombro, a cabeceira da moça emergindo dos restos da fogueira, ela toda de pé sacudindo-se como quem sai da água. (Conto \textit{Quem me deu foi a manhã}, p. 65)
\item E o tempo teve pena. Girou rápido como vento ao seu redor, como vento acariciou-lhe os cabelos e o rosto, como vento agarrou as arestas das suas roupas e do seu corpo, lambeu, soprou, desfazendo em areia aquilo que havia sido pedra, levando a areia, os grãos todos que haviam sido aquele guerreiro, para juntá-los aos infinitos grãos do seu passado. (\textit{Como cantam as pedras}, p. 36.)
\end{enumerate}

Sobre a representação da morte nos contos citados, assinale a alternativa INCORRETA:
\begin{enumerate}
\item Em \textit{Como cantam as pedras}, o corpo sólido do guerreiro é dissolvido e integrado ao infinito, permitindo o entendimento da morte como uma força natural e de contiguidade.
\item Em \textit{Na neve, os caçadores}, a transformação do caçador possibilita o entendimento da morte como metamorfose, reconstituindo a unidade familiar que ele destruíra.
\item Em \textit{Rosas na cabeceira}, a morte é apresentada a partir de uma profecia, revelando a permanência da vida nas rosas entalhadas pelo marido na cabeceira da cama.
\item Em \textit{Quem me deu foi a manhã}, a morte se materializa a partir da simbologia da salamandra e da fogueira, permitindo que a moça transponha a morte como forma de resistência.
\end{enumerate}
